\section{Related work}
\label{sec:related}

\paragraph{Calibration error in adaptive testing.} Treating estimated item parameters as
true inflates measurement precision and biases ability estimates, and adaptive selection
amplifies the problem because it seeks the items that look most informative
\citep{vanderlinden2000capitalization,olea2012capitalization,patton2013influence}. Remedies
propagate parameter uncertainty into scale scores \citep{zhang2011investigating,yang2012characterizing},
replace the item response function by its expectation over the parameter posterior---the
expected response function \citep{mislevy1994erf}, which we use---or treat all parameters in
a fully Bayesian way. \citet{vanderlinden2020fast} sample the joint posterior of ability and
item parameters during the test from stored calibration draws, and \citet{fink2025accounting}
found this approach, with selection by Bayesian maximum information, the most accurate
correction in simulated 2PL tests with 100--500 calibration responses per item. We do not
claim to be the first to integrate calibration uncertainty into adaptive testing:
Lemma~\ref{lem:exact} shows that our calibration-marginal model is the fully Bayesian model
with the calibrators' abilities and the hyperprior held at point estimates, and study~2
compares the two directly.

\paragraph{Item selection.} Classical CAT maximizes Fisher information at the current
ability estimate \citep{lord1980applications}; Bayesian variants weight it by the ability
posterior or use Kullback--Leibler information
\citep{owen1975bayesian,vanderlinden1998bayesian,chang1996global,veerkamp1997some}, and the
criteria differ mainly early in a test \citep{chen2000comparison}. Mutual-information
selection appears in classification and multidimensional CAT
\citep{weissman2007mutual,wang2011item}, as Lindley's information in experimental design
\citep{lindley1956measure,chaloner1995bayesian,rainforth2024modern} and as BALD in active
learning \citep{houlsby2011bald}, whose motivation---entropy conflates epistemic and
aleatoric uncertainty---is our Proposition~\ref{prop:decomp}; uncertainty sampling
\citep{lewis1994sequential} and expected error reduction \citep{roy2001toward} are the
counterparts of our entropy and score-targeted criteria.

\paragraph{IRT and reduced evaluation for LLMs.} IRT has been used to build evaluation
scales \citep{lalor2016building}, to weight leaderboard examples \citep{rodriguez2021evaluation},
to select anchors and compressed benchmarks
\citep{vivek2024anchor,polo2024tinybenchmarks,kipnis2025metabench} and to amortize
calibration with difficulty predictors \citep{truong2025reliable}. tinyBenchmarks choose about
100 anchors by clustering item embeddings from hundreds of models and combine them with IRT
predictions (gp-IRT); ATLAS \citep{li2026atlas} and Fluid Benchmarking
\citep{hofmann2025fluid} run Fisher-information CAT calibrated on thousands and on about one
hundred models; MINCE \citep{das2026mince} uses a few calibration models only to size a
random subset. Further work addresses continuous scores, prompt sensitivity, dataset linking,
neural and medical item banks and Laplace-approximated neural IRT
\citep{balkir2026confident,jo2025measure,habba2026growing,zhou2026lost,zheng2026medcat,mandujano2026laplace},
argues for psychometric testing of AI systems \citep{zhuang2025position}, finds item-level IRT
inferences unreliable with few calibration models \citep{jiang2026trust}---the regime of this
paper---and quantifies error bars for LLM evaluation \citep{madaan2024quantifying,miller2024adding}.

\paragraph{What is new.} The ingredients are classical. We contribute their systematic study
in the regime that LLM evaluation faces (panels of 4 to 64 systems, eight real benchmarks from
five domains, independent targets), theory for this regime, comparisons under frozen analysis
plans with the competitors that matter there, and an account of uncertainty quantification,
recalibration and cost.
